% ID: SEL-D0074-OB-E01
% Tanda S074.
Sean $x$ e $y$ los números diarios de píldoras P y Q. Dividiendo por 10 los límites de hierro y vitamina B, obtenemos
\[
 4x+y\le24,\qquad x+2y\le20,\qquad x,y\ge0.
\]
Maximizamos $C=6x+8y$ en céntimos. El conjunto factible es un polígono cerrado y acotado. En los ejes se obtienen $(0,0)$, $(6,0)$ y $(0,10)$. Para hallar el vértice restante, sustituimos $y=24-4x$ en $x+2y=20$, lo que da $x=4$, $y=8$.
\[
 \begin{array}{c|rrrr}
 (x,y)&(0,0)&(6,0)&(4,8)&(0,10)\\
 C\text{ (céntimos)}&0&36&88&80
 \end{array}
\]
El máximo del modelo es $\boxed{88\text{ céntimos al día}}$, con 4 píldoras P y 8 Q. El óptimo ya tiene coordenadas enteras y consume exactamente 240 mg de hierro y 200 mg de vitamina B.
